\documentclass[10pt,a4paper]{amsart}
\usepackage{amsmath,amssymb,amsthm,graphicx}
\usepackage{verbatim}
\newenvironment{reference}{\comment}{\endcomment}
%\newenvironment{reference}{}{}
\newenvironment{slogan}{\comment}{\endcomment}
\newenvironment{history}{\comment}{\endcomment}

% For commutative diagrams we use Xy-pic
\usepackage[all]{xy}

% We use 2cell for 2-commutative diagrams.
\xyoption{2cell}
% Theorem environments.
%
\theoremstyle{plain}
\newtheorem{theorem}[subsection]{Theorem}
\newtheorem{proposition}[subsection]{Proposition}
\newtheorem{lemma}[subsection]{Lemma}

\theoremstyle{definition}
\newtheorem{definition}[subsection]{Definition}
\newtheorem{example}[subsection]{Example}
\newtheorem{exercise}[subsection]{Exercise}
\newtheorem{situation}[subsection]{Situation}

\theoremstyle{remark}
\newtheorem{remark}[subsection]{Remark}
\newtheorem{remarks}[subsection]{Remarks}

\numberwithin{equation}{subsection}

% Macros
%
\def\lim{\mathop{\mathrm{lim}}\nolimits}
\def\colim{\mathop{\mathrm{colim}}\nolimits}
\def\Spec{\mathop{\mathrm{Spec}}}
\def\Hom{\mathop{\mathrm{Hom}}\nolimits}
\def\Ext{\mathop{\mathrm{Ext}}\nolimits}
\def\SheafHom{\mathop{\mathcal{H}\!\mathit{om}}\nolimits}
\def\SheafExt{\mathop{\mathcal{E}\!\mathit{xt}}\nolimits}
\def\Sch{\mathit{Sch}}
\def\Mor{\mathop{\mathrm{Mor}}\nolimits}
\def\Ob{\mathop{\mathrm{Ob}}\nolimits}
\def\Sh{\mathop{\mathit{Sh}}\nolimits}
\def\NL{\mathop{N\!L}\nolimits}
\def\CH{\mathop{\mathrm{CH}}\nolimits}
\def\proetale{{pro\text{-}\acute{e}tale}}
\def\etale{{\acute{e}tale}}
\def\QCoh{\mathit{QCoh}}
\def\Ker{\mathop{\mathrm{Ker}}}
\def\Im{\mathop{\mathrm{Im}}}
\def\Coker{\mathop{\mathrm{Coker}}}
\def\Coim{\mathop{\mathrm{Coim}}}

% Boxtimes
%
\DeclareMathSymbol{\boxtimes}{\mathbin}{AMSa}{"02}

%
% Macros for moduli stacks/spaces
%
\def\QCohstack{\mathcal{QC}\!\mathit{oh}}
\def\Cohstack{\mathcal{C}\!\mathit{oh}}
\def\Spacesstack{\mathcal{S}\!\mathit{paces}}
\def\Quotfunctor{\mathrm{Quot}}
\def\Hilbfunctor{\mathrm{Hilb}}
\def\Curvesstack{\mathcal{C}\!\mathit{urves}}
\def\Polarizedstack{\mathcal{P}\!\mathit{olarized}}
\def\Complexesstack{\mathcal{C}\!\mathit{omplexes}}
% \Pic is the operator that assigns to X its picard group, usage \Pic(X)
% \Picardstack_{X/B} denotes the Picard stack of X over B
% \Picardfunctor_{X/B} denotes the Picard functor of X over B
\def\Pic{\mathop{\mathrm{Pic}}\nolimits}
\def\Picardstack{\mathcal{P}\!\mathit{ic}}
\def\Picardfunctor{\mathrm{Pic}}
\def\Deformationcategory{\mathcal{D}\!\mathit{ef}}
\usepackage{cleveref}
\crefname{originallemma}{original lemma}{original lemmas}
\Crefname{originallemma}{Original lemma}{Original lemmas}
\makeatletter
\newlabel{categories-lemma-left-localization-limits}{{05Q2}{source}{}{sourcepointer-1}{}}
\newlabel{categories-lemma-left-localization-limits@cref}{{[originallemma][0][]05Q2}{[1][0][]source}{}{}{}}
\makeatother
\title{The localization of an abelian category at a multiplicative system is abelian}
\author{The Stacks Project Authors}
\date{stacks-project:homology.tex:local-snapshot-20261012; GNU FDL 1.2 or later}
\begin{document}
\maketitle
\noindent\textit{Editorial note.} All delivered mathematical text is the author's own, copied from the pinned Stacks Project chapter file homology.tex; no statement or proof step was rewritten, reordered or removed.\par\medskip
\noindent\textit{Editorial note.} The proof is the author's complete argument: part (1) is written out in full, the sentence 'Part (2) is dual to part (1).' is the author's own routine duality remark on it, and part (3) is written out after it. Every step is inside the counted span and no sentence is a placeholder.\par\medskip
\noindent\textit{Editorial note.} Every labelled reference occurring inside the counted spans is carried in assembly.external\_references as a non-load-bearing navigation pointer (same-chapter ones at their definition site, cross-chapter ones at the invocation inside the counted span), so no load-bearing mathematics is replaced by a stub. No bibliography citation occurs inside the counted spans.\par\medskip
\noindent\textit{Editorial note.} The same-name check for this card scanned every staged Stacks chapter of this repository -- all sources/stacks-*/pinned-original-source.tex files, the card's own chapter excluded -- and the labels found are recorded in dedup.same\_name\_cross\_chapter on the card.\par\medskip
\noindent\textit{Editorial note.} The preamble is the chapter's own native preamble excerpt already used for this source: the comment-type reference, slogan and history environments (native preamble lines 8-13), the Xy-pic setup (native preamble lines 14-18) and the native theorem declarations with the mathematical macros (native preamble lines 154-217, the actual last line of the file).\par\medskip
\noindent\textit{Allowed standard prerequisites (editorial).}\begin{itemize}
\item Localization and limits. Categories, Lemma categories-lemma-left-localization-limits. Supplies that the coequalizer of the zero map and Q(f) is represented by Q applied to the coequalizer of the zero map and f, which is what produces cokernels and their compatibility with the localization functor.
\item Fraction representations in a left localization. Categories, the description of a morphism of the localization as s\textasciicircum{}\{-1\} f with a common denominator. Supplies the representation used at the start of the proof to reduce statements about an arbitrary morphism of the localization to statements about an image of a morphism of the original category.
\item Abelian categories through coimage and image. Same chapter, the definition of an abelian category by the coincidence of coimage and image. Supplies the remaining axiom that part (3) has to verify.
\end{itemize}\medskip
\ifdefined\hypertarget\hypertarget{sourcepointer-1}{}\fi
\noindent\textit{Source reference.} Original lemma 05Q2: sources/stacks-categories/pinned-original-source.tex, source lines 4648--4648. Navigation pointer invoked at proof line 2151; the counted span contains the pointer only and the cited result is not reproduced inside it, so the pointer is registered as a non-load-bearing navigation stub in the established form for Stacks material. No bibliography citation is involved.\par\medskip
\begin{lemma}
\label{lemma-localization-abelian}
Let $\mathcal{A}$ be an abelian category.
\begin{enumerate}
\item If $S$ is a left multiplicative system, then
the category $S^{-1}\mathcal{A}$ has cokernels and the functor
$Q : \mathcal{A} \to S^{-1}\mathcal{A}$ commutes with them.
\item If $S$ is a right multiplicative system, then
the category $S^{-1}\mathcal{A}$ has kernels and the functor
$Q : \mathcal{A} \to S^{-1}\mathcal{A}$ commutes with them.
\item If $S$ is a multiplicative system, then the category
$S^{-1}\mathcal{A}$ is abelian and the functor
$Q : \mathcal{A} \to S^{-1}\mathcal{A}$ is exact.
\end{enumerate}
\end{lemma}

\begin{proof}
Assume $S$ is a left multiplicative system. Let $a : X \to Y$ be a morphism
of $S^{-1}\mathcal{A}$. Then $a = s^{-1}f$ for some $s : Y \to Y'$
in $S$ and $f : X \to Y'$. Since $Q(s)$ is an isomorphism we see that
the existence of $\Coker(a : X \to Y)$ is equivalent to the existence
of $\Coker(Q(f) : X \to Y')$. Since $\Coker(Q(f))$ is the
coequalizer of $0$ and $Q(f)$ we see that $\Coker(Q(f))$ is
represented by $Q(\Coker(f))$ by
Categories, Lemma \ref{categories-lemma-left-localization-limits}.
This proves (1).

\medskip\noindent
Part (2) is dual to part (1).

\medskip\noindent
If $S$ is a multiplicative system, then $S$ is both a left and a right
multiplicative system. Thus we see that $S^{-1}\mathcal{A}$ has
kernels and cokernels and $Q$ commutes with kernels and cokernels.
To finish the proof of (3) we have to show that $\Coim = \Im$ in
$S^{-1}\mathcal{A}$. Again using that any arrow in $S^{-1}\mathcal{A}$
is isomorphic to an arrow $Q(f)$ we see that the result follows
from the result for $\mathcal{A}$.
\end{proof}

\end{document}
